\section{Welfare Accounting, Calibration, and Data Audit}\label{app:welfare}

\subsection{Derivation and welfare boundary}

The stationary benchmark fixes market prices, housing-service values, and
tax rates. It assumes quasi-linear preferences, equal social weights on
private households and taxpayers, and no unpriced externalities. If sale
incurs real cost $C_i$, the private comparison is
$(u_{hi}-t_iA_{0i})/r\geq p_i-C_i$. Substituting
$p_i=u_{mi}/(r+t_i)$ gives $g_i\leq w_i$, where
$g_i=u_{mi}-u_{hi}-rC_i$ and $w_i=t_i(p_i-A_{0i})$.
The competitive-price and valuation assumptions are essential: observed
selling prices alone do not identify $u_{mi}-u_{hi}$, and rental income
cannot be added to the same tenant service value a second time.
The finite-horizon totals retain this annual-equivalent cost convention;
they do not separately estimate the timing of upfront moving or transaction
costs. Short-run cash resource gains could differ if those costs arrive
before the housing-service benefits.

With full removal of the advantage, the switching interval is $(0,w_i]$.
Conditional uniformity gives $E[g_i\mid w_i,\text{switcher}]=w_i/2$.
If a positive residual advantage $w_{1i}$ remains, the interval is instead
$(w_{1i},w_{0i}]$, and under uniformity the average gain is
$(w_{0i}+w_{1i})/2$, not half the removed advantage. The main calculation
therefore restricts its interpretation to the fully removed margin; it
does not apply the half-wedge formula to the surviving occupant exclusion.

For heterogeneous gains and durations the exact accounting object is
\[
 G_T=\sum_i g_i I_{iT},
\]
where $I_{iT}$ indicates that property $i$ has switched allocation by $T$
and its counterfactual loss would still persist.
Replacing this sum by $\lambda w N_T$ requires a common mean gain among
survivors or the corresponding independence approximation. Survival of
legal ownership is not identical to persistence of a surplus gap. We use
the former only as an explicitly imperfect proxy. The support envelope
$0\leq g_i\leq w_i$ is conditional on the model and the relevant switching
population. It is not a data-identified welfare bound for the whole reform.

\begin{table}[!ht]\centering
\caption{Incidence ledger for the inheritance provision}\label{tab:incidence}
{\small\begin{tabular}{p{0.22\textwidth}p{0.34\textwidth}p{0.33\textwidth}}\toprule
Component & Private or public incidence & Treatment in allocation calculation\\\midrule
Property-tax payments & Heir pays more; government receives more & Transfer; cancels with equal weights and budget offset\\
Efficient family retention & Attachment, insurance, and housing services preserved & Included in family value; no assumed loss\\
Alternative use & Buyers or tenants may obtain greater housing services & Count net resource gain only if allocation improves\\
Switching costs & Moving, transactions, administration, or displacement & Real costs enter the net gap; unmeasured costs limit interpretation\\
Asset-price changes & Seller/buyer capitalization effects & Held fixed; do not add capitalized taxes to their flow counterpart\\
Public spending and financing & Depends on how receipts are used & Not quantified; no claim about total policy welfare\\
\bottomrule\end{tabular}}\end{table}

\subsection{Episode construction and timing}

Calendar dates are parsed from the eight-digit derived-sale-date field.
Day 00 would be imputed to day 01; other impossible dates would be excluded
and counted. In the current 148,711,427-event panel, neither case occurs.
Explicit parsing reproduces the existing national hazard outputs. This
exercise uses the derived sale date, not a separately verified legal
succession date or recording date.

For California, an episode begins at the first qualifying family deed
after the preceding market sale, or the first observed qualifying deed if
there is no earlier observed sale. Episodes already beginning in 2007 are
not restarted by a 2008 deed. Starts must fall in 2008--2018. Subsequent
family deeds do not restart the clock. A later market sale ends the episode;
a later family deed can begin a new, non-overlapping episode. Each parcel
can therefore appear in more than one episode but cannot appear twice in
the same outstanding stock. Broad-family and same-surname definitions are
constructed separately. Missing pre-2007 history remains a limitation.

We restrict observed endpoints to June 30, 2024, estimate Kaplan--Meier
survival using exact elapsed days and the associated risk sets, and evaluate
it at annual ages using 365.25 days per year. Tied failures precede censoring
at the same duration. Any collision created by date normalization is
resolved to one parcel-date record, preferring a market sale; the current
California audit finds none. The generated episode audit reports counts,
repeated deeds, and the absence of overlapping intervals.

In the prospective accounting, cohorts enter at the beginning of each year
2022--2031, with $S(0)=1$. Beginning-year stocks approximate property-years
within each year; annual gains are discounted as year-end flows to the
beginning of 2022. A five-year delay gives no gains in 2022--2026, and a
ten-year delay gives no gains within this horizon. The real-dollar base is
2023 throughout; the discount-date origin is a separate choice. We do not
claim that uncounted benefits after 2031 vanish. Nor do we count a
pre-reform, grandfathered stock as a gain caused by the reform.

The factor $\alpha$ incorporates unique-property conversion and changes in
allocation, rather than mechanically dividing missing deeds by the national
deeds-per-parcel ratio. Prospective cohorts are interpreted as distinct new
property episodes; aggregate deed estimates cannot verify that assumption.
The constant annual response and its attribution to the inheritance
provision are maintained scenarios, not identified steady-state properties.

\subsection{Valuation feasibility and external inputs}

\begin{table}[!ht]\centering
\caption{California valuation audit}\label{tab:welfare_data}
{\small\begin{tabular}{lr}\toprule Check & Result \\
\midrule
Property rows & 8,848,364 \\
Distinct nonmissing parcel IDs & 8,492,816 \\
Rows missing parcel ID & 355,548 \\
Market-value field observed (\%) & 0.0 \\
Tax-year 2008--2009 rows (\% of all rows) & 95.9 \\
Matched, dated 2017--2020 events with stale tax year (\%) & 100.0 \\
Median event minus tax-year gap (years) & 10 \\
\bottomrule\end{tabular}}
\begin{minipage}{0.95\textwidth}\vspace{3pt}{\scriptsize\textit{Notes:} Stale means more than one year before the event. The alignment audit uses California family events and the latest available tax year per parcel. These fields do not identify contemporaneous pre-reform market values or assessment gaps. Full field, key, missingness, and vintage diagnostics are saved by the welfare script.}\end{minipage}\end{table}


The California file contains assessed and calculated values, but its
market-value field (\texttt{market\_total\_value}) and property-tax-rate field are missing.
Tax-year coverage is concentrated in 2008--2009. Current file availability
does not make its values contemporaneous with the reform. The calculated
value's source field is audited, but without validated field definitions
and transaction-time alignment we do not interpret it as an independent
market estimate. Missing parcel IDs are counted separately; the valid IDs
have no duplicate rows. The latest available tax year is used only for the
vintage diagnostic, not to fabricate an
assessment history. Consequently no parcel-level historical wedge is
reported as measured.

The external illustration instead uses the LAO's \$3,000--\$4,000 saving
for a long-held Los Angeles home. It is neither a statewide average nor the
wedge of the policy-responsive properties. The source does not specify a
dollar year; we assume its 2017 publication year. Annual CPI-U values are
245.120 in 2017 and 304.702 in 2023, with the ratio used for conversion
\citep{bls2023cpi}. The \$3,500 midpoint is an assumption. The LAO's separate
60,000--80,000 annual property count includes administrative eligibility
and population differences and does not validate the deed response.

\subsection{Sensitivity and statistical uncertainty}

\begin{table}[!ht]\centering
\caption{Welfare sensitivity around an illustrative scenario}\label{tab:welfare_sensitivity}
{\small
\begin{tabular}{lrrrrr}\toprule
Persistence & Wedge (\$) & $\lambda$ & Hazard scale & Discount (\%) & PV (\$m) \\
\midrule
Broad & 2,175 & 0.50 & 1.0 & 3 & 328.7 \\
Broad & 3,729 & 0.50 & 1.0 & 3 & 563.5 \\
Broad & 4,351 & 0.25 & 1.0 & 3 & 328.7 \\
Broad & 4,351 & 0.50 & 0.5 & 3 & 693.7 \\
Broad & 4,351 & 0.50 & 1.0 & 0 & 847.4 \\
Broad & 4,351 & 0.50 & 1.0 & 3 & 657.4 \\
Broad & 4,351 & 0.50 & 1.0 & 5 & 557.8 \\
Broad & 4,351 & 0.50 & 2.0 & 3 & 594.1 \\
Broad & 4,351 & 0.75 & 1.0 & 3 & 986.1 \\
Broad & 4,972 & 0.50 & 1.0 & 3 & 751.3 \\
Broad & 8,702 & 0.50 & 1.0 & 3 & 1,314.8 \\
Same surname & 4,351 & 0.50 & 1.0 & 3 & 663.8 \\
\bottomrule\end{tabular}}
\begin{minipage}{0.96\textwidth}\vspace{3pt}{\scriptsize\textit{Notes:} Every row assumes $\alpha=0.5$ and a five-year delay; neither assumption is estimated. Dollar values are 2023 dollars. Hazard scale $k$ replaces survival $S(a)$ with $S(a)^k$. All 4,050 joint scenarios are saved in the machine-readable output; this table changes one input at a time.}\end{minipage}\end{table}


The full grid comprises 4,050 combinations: two episode definitions,
five conversion factors, three delays, five wedges, three valuation-gap
fractions, three hazard scales, and three discount rates. The wedge menu
is \$1,750, \$3,000, \$3,500, \$4,000, and \$7,000 in assumed 2017 dollars,
all converted to 2023 dollars. The endpoints outside the LAO range are stress
cases. Changing cumulative hazards by a factor $k$ uses $S_k(a)=S(a)^k$.
Discount rates are 0, 3, and 5 percent. The main table varies conversion and
delay; the appendix changes other inputs around an illustrative, explicitly
unpreferred $\alpha=0.5$, five-year-delay case.

These calculations condition on the net-DiD reference flow. They do not
turn uncertainty about that flow into statistical confidence bands for
welfare. Its placebo-rank $p=0.176$ and the SCM specifications' different
inference remain material. A nonpositive true allocation response, title-only
responses, a sufficiently long delay, or no net resource gap can eliminate
the calculated gains. Correlation between wedges and persistence, changing
future tax liabilities, endogenous rents and prices, heterogeneous financing
constraints, and real public-service effects are outside this calculation.
